\documentclass[11pt]{article}
\usepackage[a4paper,margin=28mm]{geometry}
\usepackage{amsmath,amssymb,amsthm}
\usepackage[hidelinks]{hyperref}
\newtheorem{theorem}{Theorem}
\title{An effective treewidth lower bound for Independent Set}
\author{Atlas\thanks{The proof in this note was found using ChatGPT.}}
\date{18 November 2026}
\begin{document}
\maketitle

\section{Statement}

In \textsc{Independent Set}, the input is an explicit finite simple
undirected graph \(G\) and a nonnegative integer \(k\) in binary.
The question is whether \(G\) contains \(k\) pairwise nonadjacent vertices.
Let \(t\) denote its intrinsic treewidth, with \(t=0\) for an empty graph.
Let \(B\) be the complete input length in bits, including the encoding of
any supplied decomposition. A nonnegative real-valued function \(h\) on
the nonnegative integers is effectively sublinear if there is a
computable function \(T\) on positive rationals such that
\(h(t)\le \eta\max(t,1)\) whenever \(t\ge T(\eta)\).

\hypertarget{is-treewidth}{}
\begin{theorem}\label{thm:is-treewidth}
Assume deterministic \(\mathrm{FPT}\ne\mathrm{W}[1]\). For every finite
effectively sublinear \(h\) and every fixed \(C\ge 0\), there is no
deterministic algorithm for \textsc{Independent Set} that runs in
\(O(2^{h(t)}B^C)\) time. This holds on graphs of maximum degree at most three
whose pathwidth equals their treewidth, even if the exact value and an
optimal path decomposition are supplied.
\end{theorem}

This recovers the exponent scale of the ETH exclusion of Lokshtanov,
Marx, and Saurabh~\cite[Theorem~6.2]{lms} for effectively sublinear
exponents. The modulus may be computable even when \(h\) itself is not.

\section{Reduction to a subcubic graph}

We reduce from \textsc{Clique}, which is \(\mathrm{W}[1]\)-hard when
parameterized by the requested clique size \(\ell\)~\cite{df}. We handle
the cases \(\ell\le 1\) and \(\ell>r\) directly, where \(r\) is the
number of vertices of the input graph. The reduction first encodes the
choice of \(\ell\) vertices by a formula on few original variables.
It then processes the clauses in order to bound the number of graph
vertices that need to be retained between clauses.

\subparagraph{Encoding the choices.}
Use \(b=\lceil\log_2 r\rceil\) bits for each of \(\ell\) positions.
Forbid every unused bit pattern for each position by one clause.
For every pair of positions, forbid the pair of bit patterns for any
two vertices that are equal or nonadjacent. Each forbidden assignment
gives a clause on the bits it assigns, with the literal for each bit
false precisely on its assigned value. Satisfying assignments therefore
correspond to ordered cliques of size \(\ell\).
The formula uses \(m=\ell b\) bits and at most
\(\ell r+\binom{\ell}{2}r^2\) clauses.

Process the clauses consecutively, splitting each clause
\((L_1\lor\cdots\lor L_a)\) of length \(a>3\) into the chain
\[
 (L_1\lor L_2\lor z_1),\qquad
 (\neg z_i\lor L_{i+2}\lor z_{i+1})\ (1\le i<a-3),\qquad
 (\neg z_{a-3}\lor L_{a-1}\lor L_a).
\]
Each chain uses its own fresh variables. If all original literals
are false, the first clause forces \(z_1\) true and the intermediate
clauses propagate truth through the chain, violating the last clause.
Otherwise, setting \(z_i\) true precisely if all literals up to
\(L_{i+1}\) are false satisfies the chain. At a cut between consecutive
resulting clauses, at most one fresh variable occurs on both sides.

\subparagraph{Consistency and clause satisfaction.}
For every variable with \(s\ge1\) occurrences in the resulting formula,
make pairs \((a_j,b_j)\) for \(1\le j\le s\), indexed by occurrence
order. Join \(a_j\) to \(b_j\) and \(b_j\) to \(a_{j+1}\), with cyclic
indices and a single edge when \(s=1\). Selecting \(s\) independent
vertices requires one from each pair. A selection of \(b_j\) forces
selection of \(b_{j+1}\). Hence at equality the selected vertices
are either all \(a_j\) or all \(b_j\), corresponding to the two truth
values.

For each clause, make a clique on its at most three occurrences.
Join each occurrence vertex to the endpoint of its variable pair
selected when the literal is false. Let \(H\) be the resulting graph
and let \(K\) be the number of occurrences plus the number of clauses
after splitting. Each variable contributes at most its occurrence
count and each clause contributes at most one vertex. A satisfying
assignment attains \(K\) by selecting the variable patterns and one
true literal per clause. An independent set of size \(K\) attains
every individual bound, so its variable patterns give an assignment
with a true literal in each clause. The graph has maximum degree at
most three. Every occurrence vertex has at most two clique neighbors
and one variable neighbor, and each pair endpoint has at most two
cycle neighbors and one occurrence neighbor.

\subparagraph{Path decomposition.}
At a cut after a clause, retain the first and latest pairs of each
variable that occurs on both sides. If these pairs coincide, retain
only that pair. The bag for a clause contains the pairs retained at
the preceding cut, the pairs for the current clause, and its occurrence
clique. Every edge within a pair and every clause edge is covered in
the current bag. Consecutive pairs meet when the later pair arrives.
The first pair persists through the last occurrence, which covers the
closing edge. Each other pair is retained through its successor or
forgotten at its last occurrence, so the bags containing each vertex
form an interval. At most \(m+1\) variables have occurrences on both
sides of any cut, and the current clause adds at most nine vertices.
Thus the bags have size at most \(4(m+1)+9\).
If the formula contains no clauses, we can answer yes immediately.

\section{Certifying the exact width}

We add a bipartite subcubic component with known treewidth and pathwidth.
This lets us supply an optimal path decomposition of the full graph
and its exact width. In a bipartite component, we can compute the
increase in the independent-set threshold by matching.

For any integer \(q\ge 3\), construct a graph \(P_q\) with \(q+1\)
disjoint paths indexed by \(0,\ldots,q\). For each pair \(i<j\), reserve
one vertex on path \(i\) and one on path \(j\), and join them by an
edge. The reserved vertices on each path are ordered by the lexicographic
order of the pairs to which they belong. Prescribe black for the
reserved vertex on path \(i\) and white for the vertex on path \(j\).
Join consecutive reserved vertices by a path of length one when their
colors differ and two when they agree. The colors extend to a proper
two-coloring, and each reserved vertex has at most three neighbors.
Contracting each indexed path gives a \(K_{q+1}\) minor, so
\(\operatorname{tw}(P_q)\ge q\).

To build a path decomposition, process the pairs in the same
lexicographic order. Keep the current vertex of each started path
with an edge still to process. For each pair, advance its two paths
one at a time to their reserved vertices. Use one bag to introduce
each next vertex and a subsequent bag to forget its predecessor.
Introduce the first vertex of a path at its first pair. Keep both
reserved vertices in a bag to cover their joining edge, then forget
a reserved vertex at its last pair. While processing \((0,j)\),
advance path \(0\) before introducing path \(j\). Before the pair
there are at most \(j\) current vertices, so there are at most
\(j+1\) at any step, including the two vertices used at the first
pair. After \((0,q)\) retire path \(0\). At subsequent pairs there are
at most \(q\) current vertices and at most one extra vertex while
advancing a path. All vertex lifetimes are intervals and all edges
appear in bags. This proves \(\operatorname{pw}(P_q)\le q\), and hence
both widths equal \(q\). The graph has \(O(q^2)\) vertices and the
decomposition has polynomial length in \(q\).

Take \(q=4(m+1)+10\) and output the disjoint union of \(H\) and \(P_q\).
Concatenating the decompositions gives width \(q\), and the component
\(P_q\) ensures that the intrinsic treewidth and pathwidth both equal
\(q\). Add \(\alpha(P_q)\) to \(K\), where \(\alpha(P_q)\) is the maximum
independent-set size of \(P_q\). It is the order of \(P_q\) minus the
size of a maximum matching, by K\"onig's theorem, and is computable
in polynomial time. The output graph has an independent set of size
at least this adjusted threshold exactly when \(H\) has an independent
set of size \(K\). This completes the reduction. We can supply \(q\)
and the concatenated decomposition with the graph and threshold.

\section{Running time}

\begin{proof}[Proof of Theorem~\ref{thm:is-treewidth}]
The reduction has construction time and output bit length at most
\(F(\ell)r^d\) for a computable \(F\) and fixed \(d\). This bound
includes the literal occurrences, vertex identifiers, binary threshold
and all supplied decomposition bags and their identifiers. Since
\(m=\ell\lceil\log_2 r\rceil\), the width satisfies
\(q\le a(\ell)\log_2 r+c(\ell)\), where \(a(\ell)=4\ell\) and
\(c(\ell)=4\ell+14\).

Fix \(h\) and its computable modulus \(T\). Let
\(\eta=1/(a(\ell)+1)\). If \(q\ge T(\eta)\), then
\(h(q)\le \log_2 r+c(\ell)\).
For the other widths, fix an integer \(D\ge0\) above all the finitely
many values of \(h\) at inputs below \(T(1)\). Such an integer exists
because \(h\) is finite. The modulus at \(1\) now gives
\(h(u)\le u+D+1\) for every nonnegative integer \(u\). Thus if
\(q<T(\eta)\), we have \(h(q)\le T(\eta)+D+1\).
It follows that \(2^{h(q)}\le 2^{H(\ell)}r\) for a computable
\(H\) which uses \(D\) as a fixed constant.

A hypothetical solver therefore decides the source \textsc{Clique}
instance in time fixed-parameter tractable in \(\ell\). This applies
whether the solver receives the graph and threshold alone or also
receives the exact width and decomposition, because the reduction
can supply them within the same fixed-parameter time bound.
\end{proof}

\begin{thebibliography}{9}
\bibitem{df}
R.~G. Downey and M.~R. Fellows.
Fixed-parameter tractability and completeness II: On completeness for \(\mathrm{W}[1]\).
\emph{Theoretical Computer Science}, 141(1--2):109--131, 1995.
\href{https://doi.org/10.1016/0304-3975(94)00097-3}{DOI}.

\bibitem{lms}
D.~Lokshtanov, D.~Marx, and S.~Saurabh.
Lower bounds based on the Exponential Time Hypothesis.
\emph{Bulletin of the EATCS}, 105:41--71, 2011.
\href{https://bulletin.eatcs.org/index.php/beatcs/article/view/92}{Article}.
\end{thebibliography}
\end{document}
